\page{start}{01}{Вашият Voider}
Voider свързва поддържани Ethernet телефони чрез сдвоени уреди. Говорите по телефона. Настройвате връзката от екрана на Voider с четири бутона.
\heading{Какво е нужно}
От всяка страна: Voider с поддържан екран и бутони, подготвена microSD карта, поддържан телефон с фабрични настройки, подходящи захранвания и Ethernet кабели.

На различни места всеки Voider се нуждае от интернет с кабел през рутер. Сдвоените уреди работят и в една локална мрежа без интернет.

За сдвояване: USB памет, която може да бъде изтрита. Пазете отделно кода за проверка на всеки Voider. Това не е списък на съдържанието на комплекта.
\heading{Следваща стъпка}
\entry{wire}{Свързване на кабелите}
\entry{install}{Първа инсталация}
\entry{share}{Сдвояване на два уреда}
\entry{call}{Разговори}
\entry{help}{При затруднения}
\entry{backup}{Архив и връщане}
\note{Издание за преглед}{За кандидат 68881a782994, 20 септември 2026 г. Прочетете отворените въпроси на стр. \pageref{bg:security}.}

\page{wire}{02}{Свържете кабелите}
\wiring{Рутер / локална мрежа}{Мрежов порт (eth0)}{Порт за телефон (eth1)}{Телефон: мрежов порт}{Схема на връзките, а не на местата на портовете.}
\begin{enumerate}
\item Свържете кабелите при изключено захранване.
\item Телефонният мрежов порт свържете към телефонния порт на Voider, не към рутера.
\item Мрежовия порт на Voider свържете към LAN на рутера.
\item Включете предвидените захранвания. Повторете от другата страна.
\end{enumerate}
Виждате настройката или \UI{DEVICE CHECK}. След проверките \UI{NET} показва \UI{PHONE}: \UI{CONNECTED} за достъпен телефон. Разговорът още трябва да се провери.
\note{Първо уточнете хардуера}{Надписите, местата на портовете, захранването и одобрените телефони не са документирани. Уточнете ги преди свързване. Не приемайте, че Ethernet подава ток.}

\page{controls}{03}{Бутони и език}
Четири бутона, четири полета долу на екрана: съответстват отляво надясно. Надписите се сменят. Празното поле няма действие.

\textbf{Натиснете:} натиснете и отпуснете. \textbf{Задръжте:} натиснете за около секунда, после отпуснете. Две бели черти над надписа означават задържане. Прочетете новия екран преди следващото натискане.
\shot{language}{Изображение от истинския софтуер на екрана с примерни данни.}
\begin{enumerate}
\item От началния екран: \UI{NET}, после \UI{SYSTEM}, после \UI{LANGUAGE}.
\item С \UI{NEXT} изберете Български. Натиснете \UI{TRY}. Езикът се сменя до изключване.
\item За постоянен запис задръжте \UI{SAVE}. Запишете и потвърдете новия код (стр. \pageref{bg:check}).
\end{enumerate}
Настройката започва на английски: \textbf{Language}, \textbf{Next}, \textbf{Try}, после \UI{BACK}. Запазете след началните проверки.

\page{install}{04}{Първа инсталация}
Само при \UI{INSTALL VOIDER}. При \UI{DEVICE CHECK} продължете на следващата страница.
\shot{factory-READY}{За начало задръжте втория бутон.}
\begin{enumerate}
\item При нужда изберете език (стр. \pageref{bg:controls}). Задръжте \UI{INSTALL}.
\item Изчакайте. Етапът и процентът се сменят. Оставете захранването включено. Всичко е на картата; не е нужно изтегляне от интернет.
\item При \UI{INSTALL COMPLETE} запишете кода и за кой уред е. Това е първото копие за сравнение.
\item Натиснете \UI{REBOOT} веднъж. Екранът потъмнява при рестарта. Изчакайте \UI{DEVICE CHECK}.
\end{enumerate}
Инсталацията създава уникална самоличност. Повторете на другия Voider; пазете кода отделно. Няма автоматичен рестарт.

При \UI{INSTALL FAILED} запишете съобщението. След поправка задръжте \UI{RETRY}, или \UI{POWER} за изключване.

\page{check}{05}{Проверка при всеки старт}
\begin{enumerate}
\item Изчакайте \UI{DEVICE CHECK} да покаже \UI{PASSED}. Инсталираната система е проверена. Натиснете \UI{NEXT}.
\item В \UI{YOUR CHECK} сравнете кода със записа за \emph{този} Voider. \UI{FULL} показва целия отпечатък; \UI{BACK} връща назад.
\item Ако съвпада, натиснете \UI{MATCH}. Виждате \UI{PASSED}. Натиснете \UI{HOME}. Връзките вече могат да започнат.
\end{enumerate}
\shot{integrity_state}{Само пример. Не използвайте този код като свой запис за сравнение.}
\textbf{Разлика или липсващ запис?} Не гадайте. При неочаквана разлика задръжте \UI{DIFFER}, после \UI{POWER}. Потърсете причината. Неуспешната проверка на уреда също спира работата.

\textbf{След Ваша промяна:} при \UI{STATE UPDATED} запишете новия код. Проверете записа си и натиснете \UI{I RECORDED IT}. Връщате се директно към предишния екран. \UI{MATCH}/\UI{DIFFER} служи за сравнение със записания код при стартиране. Сдвояване, запис и изтриване могат да го сменят. Пазете стария запис с причината. Нормален рестарт не го сменя.

\page{share}{06}{Сдвояване: първи уред}
Сдвояването дава на двата уреда данни за връзката. Нека ги наречем A и B. Това не са надписи на екрана. Едно прехвърляне стига за разговори в двете посоки.
\transfer{A: \UI{SHARE}}{Пренесете USB}{B: \UI{ADD}}
\begin{enumerate}
\item Поставете една USB памет и извадете другите. На A натиснете \UI{LINKS}, после \UI{SHARE}. Voider избира паметта веднага. Ако липсва, поставете я и натиснете \UI{RETRY}.
\item При \UI{ERASE USB} проверете показания модел. Задръжте \UI{ERASE} само ако всички файлове могат да бъдат загубени. \UI{CANCEL} отменя без изтриване.
\item Изчакайте \UI{STATE UPDATED}. Запишете и потвърдете новия код на A (стр. \pageref{bg:check}).
\item При \UI{SAFE TO REMOVE} извадете паметта и натиснете \UI{DONE}. Занесете я направо на B.
\end{enumerate}
\note{Пазете паметта}{Тя съдържа тайни данни за връзката. Voider не ги шифрова. Не я давайте на други хора. За всяка друга връзка започвайте ново \UI{SHARE}.}

\page{add}{07}{Сдвояване: втори уред}
\begin{enumerate}
\item Поставете USB паметта от A. На B натиснете \UI{LINKS}, после \UI{ADD}. Voider избира паметта веднага. Ако липсва, поставете я и натиснете \UI{RETRY}. Появява се \UI{VALID VOIDER KEY}. Няма отпечатък за сдвояването. Файлът вече е прочетен и проверен без промяна на паметта.
\item Задръжте \UI{ADD} отново за внасяне. Паметта се чете, не се изтрива. Изчакайте \UI{STATE UPDATED}.
\item Запишете и потвърдете кода на B.\par При \UI{SAFE TO REMOVE} извадете паметта, натиснете \UI{DONE} и я приберете.
\item Оставете двата уреда включени. Отворете връзката и изчакайте \UI{CONNECTED}. После пробвайте разговор.
\end{enumerate}
\note{Две различни проверки}{Началните кодове на A и B са за различни уреди. Не е нужно да са еднакви. Сравнявайте всеки само с неговия отделен запис.}
\heading{Липсва отпечатък на екрана}
Софтуерът пази отделен отпечатък на USB файла. Този кандидат няма екран за сравняването му между A и B. \UI{VALID VOIDER KEY} не проверява самоличност на човек. Ако Ви е нужно това сравнение, спрете тук до решаване на липсата. Началните кодове не го заменят.

При грешка вижте стр. \pageref{bg:help}. Не споделяйте отново само защото връзката още се изчаква.

\page{call}{08}{Обаждане и отговор}
\begin{enumerate}
\item Натиснете \UI{LINKS}, после \UI{OPEN}. Изберете с \UI{PREV} или \UI{NEXT}.
\item Изчакайте \UI{CONNECTED}. Прочетете адреса в \UI{DIAL}.
\item Във функцията \textbf{Direct IP Call} на телефона въведете този адрес и започнете обаждането.
\item Другият телефон трябва да звъни. Вдигнете слушалката за отговор. Проверете звука в двете посоки. Затворете за край.
\end{enumerate}
\shot{contact}{Пример: тук се набира 10.1.2.1. Използвайте своя показан адрес.}
Бутоните на Voider не набират. Всеки може да се обади. Променящ се входящ адрес 172.16.x.x не е за набиране. За обратно обаждане използвайте екрана на Voider.
\note{Липсват телефонни стъпки}{Одобреният модел, менюто Direct IP Call и точните бутони за набиране трябва да се потвърдят преди доставка на клиенти.}

\page{manage}{09}{Управление на връзките}
\textbf{Друг човек:} повторете стр. \pageref{bg:share}–\pageref{bg:add} с ново \UI{SHARE} на A. И попълненият списък има \UI{SHARE} и \UI{ADD}. Не сдвоявайте два пъти за разговори в двете посоки.

\UI{CLIENT} и \UI{SERVER} са местни роли, не имена на хора. Споделянето създава запис Клиент на A; добавянето --- запис Сървър на B. Запишете кой е към всеки запис. Екранът не предлага преименуване.

\textbf{Изтриване на една връзка:}
\begin{enumerate}
\item \UI{LINKS} \arr \UI{OPEN}; изберете с \UI{PREV}/\UI{NEXT}.
\item \UI{MORE} \arr \UI{REMOVE}. Проверете името, после задръжте \UI{REMOVE}.
\item Запишете и потвърдете новия код. Записът изчезва само тук. Другият списък остава. Този уред отказва повторно внасяне на премахнатия USB файл.
\end{enumerate}

\textbf{Повторно свързване:} Отворете връзката, изберете \UI{MORE} \arr \UI{RETRY}, после задръжте \UI{RETRY}. Сдвояването се запазва; другите връзки продължават да работят.

\textbf{Нулиране:} \UI{NET} \arr \UI{SYSTEM} \arr \UI{MORE} \arr \UI{RESET}, после \UI{NEXT} и задръжте \UI{RESET}. Изтрива всички връзки и SSH ключа за достъп, изключва SSH и създава нови самоличности. Нужно е ново сдвояване. Някои настройки, включително езикът, остават. Запишете новия код. Не е решение за бавна връзка.

\page{status}{10}{Състояние на връзката}
\UI{NET} показва интернет и телефон. \UI{PHONE}: \UI{CONNECTED} означава, че Voider достига телефона.

\textbf{Броят на началния екран:} 1/2 значи, че една от две записани връзки е свързана. Клиенти и сървъри се броят отделно. Това не е брой разговори.

\textbf{\UI{CONNECTED}:} има активен път до другия Voider. \textbf{\UI{CONNECTING}:} прави се опит. \textbf{\UI{OFFLINE}:} не се показва активен път. Това не доказва, че другият Voider е изключен.

Voider обикновено опитва \UI{LAN}, после \UI{DIRECT}, после \UI{PRIVATE} (WireGuard през рутери), после \UI{TOR}. Показаният път е за тази връзка. LAN също използва WireGuard. Връзките показват \UI{LAN4}, \UI{LAN6}, \UI{DIRECT4} или \UI{DIRECT6} според избраната IP версия.

\textbf{Дайте му време.} Проверките, сверяването на часовника и стартът на Tor могат да отнемат минути. Опитът за възстановяване на Tor има прозорец от 90 секунди. Цялото свързване може да отнеме повече. Няма гарантиран срок. Оставете двата уреда включени, докато мрежата се връща. Ново сдвояване не ускорява това.

В една локална мрежа активна LAN връзка може да работи без интернет. Достъпността на телефона се проверява отделно.

\textbf{\UI{PATHS}:} обикновено оставете седемте включени. За промяна: \UI{UP}/\UI{DOWN} избира, \UI{REMOVE}/\UI{ADD} превключва. \UI{BACK}, после задържане на \UI{APPLY} записва; потвърдете кода. \UI{CANCEL} връща към редактиране. Поне един път остава включен. Ограниченията могат да спрат обажданията.

\page{help}{11}{При затруднения}
\begin{enumerate}
\item Проверете двата екрана. Началните проверки трябва да са завършени и на двата уреда.
\item Натиснете \UI{NET}. При липсващ телефон проверете тока и кабела му към Voider. При липсваща мрежа проверете кабела към рутера и самия рутер.
\item Оставете двата уреда включени за няколко минути. Още няма връзка? В \UI{NETWORK} натиснете \UI{RETRY}, после го задръжте на следващия екран. Това подновява проверките и може да прекъсне разговор.
\item Ако не помогне, запишете съобщението и избраната връзка. Пазете сдвояването и кодовете. Не изтривайте и не нулирайте като обичаен начин за поправка.
\end{enumerate}
\textbf{USB:} Екранът различава липсваща, изключена или сменена памет и грешки при четене, запис и проверка. Използвайте една памет. Монтирани или опасни носители се отказват. След грешка отворете действието и изберете паметта отново.

\textbf{Няма звук или звънене:} проверете адреса и двата телефона. Зелена връзка или достъпен телефон не доказва звук. Уеб управлението на телефона не е нужно.

\textbf{Нормално изключване:} начало \arr \UI{POWER}, после задръжте \UI{POWER}. Виждате \UI{POWERING OFF}. Данните се записват; екранът потъмнява при началото на изключването. Изчакайте, преди да спрете тока. Черен екран сам по себе си не доказва, че процесът е завършил.

\textbf{Нормален рестарт:} Начало \arr \UI{POWER}, после \UI{REBOOT} вдясно. На екрана за потвърждение задръжте \UI{REBOOT}. Запазените връзки остават; след стартиране завършете нормалните проверки.

\page{backup}{12}{Запазете архив}
Архивът съдържа самоличността, връзките, ключовете и настройките на този уред. Не архивира телефона или цялата системна карта. Файловете не са шифровани. Пазете паметта като ключ за уреда.
\begin{enumerate}
\item От началото: \UI{NET} \arr \UI{SYSTEM} \arr \UI{BACKUP}.
\item Поставете резервна USB памет. Задръжте \UI{BACKUP}.
\item Прочетете \UI{ERASE USB}. Всички файлове ще се изтрият. Ако това е желано, задръжте \UI{ERASE}.
\item Изчакайте \UI{SAFE TO REMOVE}. Извадете паметта. Запишете уреда, датата и текущия код. Пазете отделни архиви за A и B.
\end{enumerate}
\heading{Връщане на архив}
Поставете паметта с архива. Отворете \UI{NET} \arr \UI{SYSTEM} \arr \UI{MORE} \arr \UI{RESTOR} и натиснете \UI{CHECK}. Voider чете и проверява архива, без да променя паметта. Инсталираната система трябва да съвпада.

Задръжте \UI{RESTOR} само за да замените всички записани настройки, самоличности и връзки с архива. По-новите промени се губят. Запишете новия код и натиснете \UI{I RECORDED IT}. Извадете паметта и задръжте \UI{REBOOT}. Старите работни данни не могат да презапишат архива. След рестарта сравнете кода с \UI{MATCH} или \UI{DIFFER}.

\page{ssh}{13}{По избор: SSH достъп}
SSH е за доверен администратор, не за разговорите. Отворете \UI{NET} \arr \UI{SYSTEM} \arr \UI{MORE} \arr \UI{SSH}. Виждате \UI{MANAGEMENT ON} или \UI{MANAGEMENT OFF}.

\textbf{Включване или спиране:} натиснете \UI{ON} или \UI{OFF}, после задръжте същия бутон за потвърждение. Запишете и потвърдете новия код. Спирането забранява нов вход, но пази ключа. Текущите сесии и вътрешният пренос остават.

\textbf{Създаване или смяна на ключ:}
\begin{enumerate}
\item Изберете \UI{LOGIN KEY}. Прочетете въпроса за създаване или смяна. При смяна старият ключ спира да работи. Задръжте \UI{CONFIRM}.
\item Поставете една памет за изтриване и натиснете \UI{CHECK}. Задръжте \UI{ERASE}: всички дялове и файлове се изтриват. Достъпът се сменя едва след запис и проверка.
\item Запишете и потвърдете кода. SSH е включен. При \UI{SAFE TO REMOVE} извадете паметта и натиснете \UI{DONE}. Пазете я. На нея са \texttt{voider-admin} (частен ключ за вход) и \texttt{voider-admin.pub} (публичен). Частният ключ няма парола.
\end{enumerate}
\textbf{Окончателно премахване:} изберете \UI{REMOVE}, после задръжте \UI{REMOVE}. Ключът се отменя и SSH остава изключен. Запишете новия код. При USB грешка старият ключ и настройката за достъп остават. Voider пази само публичната част; изгубен частен ключ не може да се възстанови.

\page{security}{14}{Проверки и защита}
\textbf{Проверка на уреда:} сравнява инсталираната система със записани стойности. \textbf{Твоята проверка:} сравнявате запазеното състояние със своя запис. Нито една не доказва кой държи другия телефон.

\textbf{Сдвояването} създава споделени тайни. Отпечатъкът за сдвояване не се показва за сравнение (стр. \pageref{bg:add}). Пазете USB паметите и архивите.

\textbf{Защита в мрежата:} WireGuard защитава LAN, преките пътища и тези през рутери. По пътя през Tor шифроването се осигурява от Tor. tundup пренася удостоверен явен текст: проверява автентичността и отхвърля повторения, но сам не шифрова данните на разговора. Не приемайте, че локалният телефонен кабел е шифрован.

Собственият протокол няма независим одит. Voider не гарантира анонимност и не предотвратява анализ на трафика. Откраднати тайни за сдвояване могат да застрашат и минали сесии.

\textbf{SSH граници:} root достъп само с ключ от частни мрежи през мрежовия Ethernet порт. Без пароли, вход от телефонния порт, публичен WAN или Tor; без препращане. Подготвен ключ включва SSH при инсталация; без него е изключен. После важат двете проверки. Администраторът проверява ключа на хоста. За това и за SSH превключване при инсталация няма екран.

\mini{\textbf{Преди доставка:} уточнете портове, захранване, телефонни бутони и инсталационен SSH. Уточнете проверката при сдвояване. Остават езиков преглед, тест с начинаещи, нови инсталации и финални тестове. Tor е наблюдаван 601 секунди, не час. Образът не е обявен за готов за продажба.}
